% answers_short.tex —— 2026 年·北京卷数学（回忆版，内部候选） · 简短答案速查
% 由 tools/make_short_answers.py 自 answers.tex 的 \gksolhead/\gkanswer 提取，
% 并与 2026/2026_数学真题_北京回忆版_答案要点.md 逐题比对后生成，勿手改（改源 MD 后重跑）。

\begin{gknotice}
\textbf{使用说明}：本节只给最终答案便于快速核对，完整推导见解析版。答案为编写者独立推导，\textbf{非官方标准答案}。
\end{gknotice}

\begin{gknotice}
选择题答案中的 U 表示待核，未给出猜测答案，详情见解析版。
\end{gknotice}

\gkpart{单选题}
\begin{center}\begin{tabular}{ccccccccccc}
\toprule
题号 & 1 & 2 & 3 & 4 & 5 & 6 & 7 & 8 & 9 & 10 \\
\midrule
答案 & B & A & B & A & D & D & A & C & B & U \\
\bottomrule
\end{tabular}\end{center}

\gkpart{填空题}
\begin{itemize}[leftmargin=2.6em]
  \item[\textbf{11.}] \(0\)
  \item[\textbf{12.}] \(d=-2\)；\(a_1\in[8,10]\)，例如 \(a_1=8\)
  \item[\textbf{13.}] U/PENDING；固定回忆版候选为 \([700,4900]\)，第三方异文对应 \([700,4900)\)
  \item[\textbf{14.}] 底面积 \(\sqrt3\)，体积 \(\dfrac{2\sqrt3}{3}\)
  \item[\textbf{15.}] ①②③④（②、③的原文序号有异文；两版下四项均成立）
\end{itemize}

\gkpart{解答题}
\begin{itemize}[leftmargin=2.6em]
  \item[\textbf{16.}] (1) \(\omega=2,\ \varphi=\dfrac\pi3\)；(2) \(\left[\dfrac\pi{12}+k\pi,\dfrac{7\pi}{12}+k\pi\right]\)，\(k\in\mathbb{Z}\)
  \item[\textbf{17.}] (1) \(0.8\)；(2) 独立多项分布模型下 \(\dfrac6{625}\)；(3) \(s_{\rm 左}^2<s_{\rm 中}^2<s_{\rm 右}^2\)
  \item[\textbf{18.}] (1) 见解析版；(2) U/PENDING；若取条件①或②，夹角余弦均为 \(\dfrac23\)；条件③不能唯一确定 \(P\)
  \item[\textbf{19.}] (1) \(\dfrac{x^2}{4}+\dfrac{y^2}{3}=1\)；(2) \(k=\pm\dfrac{\sqrt2}{4}\)；\(k=-\dfrac34\) 时 \(Q\) 不存在
  \item[\textbf{20.}] (1) \(m=2,n=1\)；(2) 2 个极值点；(3) 1 个交点
  \item[\textbf{21.}] (1) \(A_1\) 否，\(A_2\) 是；(2) 最多 12 个 1；(3) \(a_{ij}=a_{i1}a_{1j}\)
\end{itemize}
